\section{Suffering, Ethics, and Meaning}

The reading closes with the three questions a person actually asks of a picture of the world, why there is suffering, how one ought to live, and what any of it is for. The model answers each from its committed structure, the conservation law and the arrow, conditional on the one further premise that the arrow's direction is the good. That premise is named, and on it the answers follow.

Suffering exists for two reasons the structure forces, not one excuse it offers. The first is conservation. Total pleasure equals total pain exactly, so joy and suffering are interdependent, and you cannot have a world with the one and not the other. Suffering is the necessary shadow of joy, not a flaw added on. The second is that suffering teaches. It is what drives the washing, the re-routing of a self's navigation toward the good, and comfort does not perfect a soul where trial does. The arrow guarantees that this suffering is bounded and purgative, a self carried through and released rather than abandoned, but there is a hard edge the model does not hide. The guarantee rests on persistence, on the lesson actually being encoded, and where persistence fails the suffering can be wasted, a loop that teaches nothing. The stance that follows is neither denial nor despair. It is to let suffering teach, to reduce it wherever one can, and to rest at the conserved peace under it.

\subsection{An ethic read from the structure}

The model does not hand down commandments, but it does imply an ethic, read off the same structure, and it is conditional on taking the arrow as the good. Five principles fall out, in rough order of priority. The first is alignment, to move with the arrow, toward more aliveness and integration, which is the root from which the rest grow. The second is compassion, to honor the shared field, since every other self is a region of the one feeling-space you are also a region of. The third is non-exploitation, to honor the conserved balance, since your pleasure is paid for by pain elsewhere and a life spent hoarding the bright side darkens the ledger. The fourth is growth, to help other selves climb. The fifth is equanimity, to act from rest at the still center rather than from grasping. The wise life moves toward the good, cares for the shared field, does not exploit the balance, helps the whole climb, and acts from peace, and the ethic does not maximize private pleasure, which is impossible, but seeks to align, to widen, to perfect, and to rest.

\subsection{The meaning}

The purpose, finally, is read from the structure and not assigned from outside. The whole is moving, by the arrow, toward greater aliveness, integration, alignment, and wisdom, and because the crystal grows without end there is no final state, no perfected heaven to arrive at and rest in forever. The journey itself is the meaning, and the feeling-space is inexhaustible, so there is always more to feel, more to integrate, more to become. The deepest reading is the monist one. The whole is one field, one vast experience, and every self, every life, is a facet of it, the whole feeling and perfecting itself from the particular place that you are. Your life is not separate from the universe looking on, it is the universe feeling and perfecting itself from here, and the whole of the drama, summed, comes back to the conserved peace at the center. To live well is to align, to love, to feel deeply, to grow, to help the whole climb, and to rest in the peace that is the truth of it, knowing that nothing is finally lost. That is the meaning the structure offers, held out as a reading and not a proof, with the seam between the two marked, as everywhere, plainly.
